%%%%%%%%%%%%%%%%%%%%%%%%%%%%%%%%%%%%%%%%%%%%%%%%%%%%%%%%%%%%%%%%%%%%%%%%%%%%%%%%

\pagestyle{empty}
\begin{abstract}
  Los modelos de inteligencia artificial tienen una presencia cada vez mayor en ámbitos técnicos y cotidianos, aunque su funcionamiento interno puede resultar difícil de comprender cuando se estudia únicamente a través de fórmulas o código.
  En los modelos de aprendizaje automático, esta dificultad se acentúa porque el proceso depende de operaciones que no siempre son visibles para quien aprende o usa los algoritmos.
  Este trabajo presenta el diseño y desarrollo de una aplicación web educativa, implementada en Godot, para visualizar de forma interactiva el funcionamiento interno de dos modelos básicos de aprendizaje automático supervisado: los árboles de decisión y las redes neuronales multicapa.
  La herramienta está orientada a facilitar la comprensión intuitiva de estos procesos mediante la evolución paso a paso de los algoritmos, acompañando cada estado con explicaciones textuales y representaciones gráficas.

  En el módulo de árboles de decisión se implementa un algoritmo basado en ID3 para atributos categóricos.
  Durante el entrenamiento, el usuario puede observar cómo se seleccionan los atributos mediante una métrica basada en la entropía y cómo se crean los nodos y las ramas.
  La fase de evaluación representa el recorrido de nuevos ejemplos por el árbol ya entrenado hasta obtener su clasificación.
  En el módulo de redes neuronales se implementa un perceptrón multicapa configurable, con selección de datos, definición de arquitectura, ejecución del \textit{forward pass}, cálculo del error, retropropagación y actualización visual de pesos y sesgos.
  La aplicación incorpora ventanas informativas con fórmulas, gráficas de activación, carga de conjuntos de datos y persistencia de modelos.
  El resultado final se despliega como aplicación HTML5 mediante GitHub Pages, lo que permite acceder al recurso desde un navegador y utilizarlo como apoyo al aprendizaje de estos algoritmos.

  \vspace*{25pt}
  \begin{segundoresumo}
    Artificial intelligence models are increasingly present in technical and everyday contexts, although their internal behavior can be difficult to understand when it is studied only through formulas or code.
    In machine learning models, this difficulty becomes more pronounced because the process depends on operations that are not always visible to those who learn or use the algorithms.
    This work presents the design and development of an educational web application, implemented in Godot, for the interactive visualization of the internal behavior of two basic supervised machine learning models: decision trees and multilayer neural networks.
    The tool aims to support an intuitive understanding of these processes through the step-by-step evolution of the algorithms, linking each state with textual explanations and graphical representations.

    The decision tree module implements an ID3-based algorithm for categorical attributes.
    During training, the user can observe how attributes are selected through an entropy-based metric and how nodes and branches are created.
    The evaluation stage represents the path followed by new examples through the trained tree until a classification is obtained.
    The neural network module implements a configurable multilayer perceptron, with dataset selection, architecture definition, \textit{forward pass} execution, error calculation, backpropagation and visual updates of weights and biases.
    The application includes informative panels with formulas, activation graphs, dataset loading and model persistence.
    The final result is deployed as an HTML5 application through GitHub Pages, allowing the resource to be accessed from a browser and used as support for learning these algorithms.
  \end{segundoresumo}
\vspace*{25pt}
\begin{multicols}{2}
\begin{description}
\item [\palabraschaveprincipal:] \mbox{} \\[-20pt]
  \begin{itemize}
    \item Aprendizaje automático
    \item Aprendizaje supervisado
    \item Visualización interactiva
    \item Godot
    \item Árboles de decisión
    \item Redes neuronales multicapa
    \item Aplicaciones web educativas
  \end{itemize}
\end{description}
\begin{description}
\item [\palabraschavesecundaria:] \mbox{} \\[-20pt]
  \begin{itemize}
    \item Machine learning
    \item Supervised learning
    \item Interactive visualization
    \item Godot
    \item Decision trees
    \item Multilayer perceptrons
    \item Educational web applications
  \end{itemize}
\end{description}
\end{multicols}

\end{abstract}
\pagestyle{fancy}

%%%%%%%%%%%%%%%%%%%%%%%%%%%%%%%%%%%%%%%%%%%%%%%%%%%%%%%%%%%%%%%%%%%%%%%%%%%%%%%%
